
\def\PUDcodigo{ECA213}

\def\PUDcodacad{04507.28}

\def\PUDdisciplina{Microcontroladores}

\def\PUDsemestre{S6}

\def\PUDhoras{80}

%NOVOS ---------------------------------------------------
\def\PUDhorasTeoria{60}
\def\PUDhorasPratica{20}

\def\PUDinterECA{ 
\hyperref[PUD:ECA208]{Lógica de Programação (04507.10)} 

\hyperref[PUD:ECA209]{Linguagem de Programação (04507.14)}
}

\def\PUDatuaisECA{---}

\def\PUDrecursos{Projetor multimídia; Quadro branco e pincel; Aulas em laboratório.}

%----------------------------------------------------------

\def\PUDcreditos{4}

\def\PUDprerequisito{ \hyperref[PUD:ECA209]{ECA209} }

\def\PUDpreacad{ \hyperref[PUD:ECA209]{Linguagem de Programação (04507.14)} }

\def\PUDobjetivos{Compreender o principio básico de funcionamento de um microcontrolador;  Conhecer e utilizar simuladores de microcontroladores;  Compreender o princípio de clock, temporização e reset de um microcontrolador;  Programar um microcontrolador utilizando Linguagem C;  Utilizar os dispositivos de interrupções, temporização e comunicação serial dosmicrocontroladores.  Projetar e analisar sistemas eletrônicos microcontrolados.}

\def\PUDementa{Introdução;  Arquitetura Interna de Microcontroladores RISC;  Estudo dos Pinos do Microcontrolador;  Clock, Ciclos de Temporização e Reset;  Conjunto de Instruções;  Sistemas de Interrupção;  Temporizadores e Contadores;  A Comunicação Serial;  Modos de baixo consumo.  Projetos Práticos.}

\def\PUDconteudo{ &  INTRODUÇÃO: Definição de microcontrolador, Tipos e arquitetura dos microcontroladores, Introdução à linguagem assembly, Revisão de Linguagem C para Microcontroladores, Apresentação de simuladores e compiladores. 
\\ 
 &  ARQUITETURA INTERNA DE MICROCONTROLADORES RISC: Arquitetura da ULA, Funções das FLAGs, Registradores de uso geral e de uso específicos, Instrução/Operando, Executando um programa passo a passo, Estudo da Memória Interna e Externa.  Apresentação das famílias PIC 18F, MSP430 e ARM. 
\\ 
 &  ESTUDO DOS PINOS DO MICROCONTROLADOR: Descrição da pinagem, Descrição das funções, Aplicações Práticas. 
\\ 
 &  CLOCK, CICLOS DE TEMPORIZAÇÃO E RESET: Geração de Clock, Tempos de Processamento, Estudo do Reset. 
\\ 
 &  CONJUNTO DE INSTRUÇÕES: Tipos de instruções, Estudo do conjunto de instruções, Exemplos básicos de sub-rotinas. 
\\ 
 &  SISTEMAS DE INTERRUPÇÃO: Estrutura da interrupção, Tipos de interrupções, registros especiais e suas programações, Aplicações Práticas. 
\\ 
 &  TEMPORIZADORES E CONTADORES: Modos de funcionamento, Registros Especiais e suas programações, Aplicações Práticas. 
\\ 
 &  A COMUNICAÇÃO SERIAL: Características básicas da comunicação serial, A interface serial, Modos de programação, Taxas de Transmissão (Baud-rate), Protocolo de Comunicação serial RS 232. 
\\ 
 &  MODOS DE BAIXO CONSUMO: Configuração.  Estratégias. 
\\ 
 &  PROJETOS PRÁTICOS: Projetos utilizando microcontroladores – Práticas de software e hardware em laboratório. }

\def\PUDmetodo{Aulas expositórias que podem ser teóricas e/ou práticas, onde as práticas no laboratório serão marcadas no decorrer da disciplina.}

\def\PUDavalia{A avaliação será feita com aplicação de provas teóricas e/ou práticas, além da possibilidade de inclusão de trabalhos e seminários no decorrer da disciplina.}

\def\PUDbibbasica{ &   \href{www.biblioteca.ifce.edu.br/index.asp?codigo_sophia=24390}{Pereira}, Fábio.  Microcontroladores PIC – Programação em C.  7º ed.  São Paulo, SP: Érica, 2007.  358p.  ISBN: 9788571949355.
\\ 
 &   \href{www.biblioteca.ifce.edu.br/index.asp?codigo_sophia=24030}{Souza}, Daniel Rodrigues de.  Desbravando o microcontrolador PIC18: recursos avançados.  São Paulo, SP: Érica, 2010.  336p.  ISBN: 9788536502632.
\\ 
 & \href{www.biblioteca.ifce.edu.br/index.asp?codigo_sophia=23899}{Sousa}, Daniel Rodrigues de.  Microcontroladores ARM7: (Philips - Família LPC213X): o poder dos 32 bits: teoria e prática.  São Paulo, SP: Érica, 2006.  278p.  ISBN: 8536501200. }
 %&   PEREIRA, Fábio.  Tecnologia ARM - Microcontroladores de 32 bits.  1a.  Ed.  Érica, 2007.  448 p.  ISBN 9788536501703. }

\def\PUDbibcomplementar{ & \href{www.biblioteca.ifce.edu.br/index.asp?codigo_sophia=24331}{Nicolosi}, Denis E. C.  Laboratório de Microcontroladores: Família 8051, Treino de instruções, hardware e software.  5º ed.  São Paulo, SP: Érica, 2006.  206p.  ISBN: 8571948712.
\\ 
 &  \href{www.biblioteca.ifce.edu.br/index.asp?codigo_sophia=24174}{Nicolosi}, Denis E. C.; Bronzeri, Rodrigo B.  Microcontrolador 8051 com linguagem C - Prático e Didático - Família AT89S8252.  2º ed.  São Paulo, SP: Érica, 2008.  222p.  ISBN: 9788536500799.
\\ 
 &  \href{www.biblioteca.ifce.edu.br/index.asp?codigo_sophia=23871}{Souza}, David José de.  Desbravando o PIC: ampliado e atualizado para PIC 16F628A.  12º ed.  São Paulo, SP: Érica, 2008.  268p.  ISBN: 9788571948679.
\\
 &  \href{www.biblioteca.ifce.edu.br/index.asp?codigo_sophia=67116}{Jucá}, Sandro.  Aplicações práticas de microcontroladores utilizando software livre: aprenda de forma prática a gravação wireless e via usb de microcontroladores através da ferramenta SanUSB.  Recife, PE: Imprima, 2017.  200p.  ISBN: 9788564778610.
\\
 &  \href{www.biblioteca.ifce.edu.br/index.asp?codigo_sophia=55618}{Gimenez}, Salvador P.  Microcontroladores 8051: teoria do Hardware e do Software: aplicações em controle digital: laboratório e simulação.  E-book.  Pearson.  272p.  ISBN: 9788587918284. }

\def\PUDautor{Geraldo Luis Bezerra Ramalho}

\def\PUDdata{2013-04-17}

\def\PUDrevisao{3}

\def\PUDrevisor{Fábio Timbó Brito}

\def\PUDdatarevisao{2019-05-15}

\PUDcorpo
